% Proposed figure inclusions and captions for the Branin motivational
% experiments (impl_plan/motivational_exp.txt). Captions describe what was
% measured, not a claimed conclusion. Paths are relative to
% experiments/eval2_branin_motivation/results/figures/.

% ============================================================
% Experiment A: where should cross-task knowledge enter BO?
% ============================================================

\begin{figure}[t]
  \centering
  \includegraphics[width=\textwidth]{figures/motivation_A_task_scaling.pdf}
  \caption{Simple regret to the numerically verified global optimum
  (mean $\pm$ s.e.m.\ over 20 held-out Branin tasks) as a function of the
  number of completed training tasks, measured immediately after
  initialization (left), after 10 BO calls (center), and after the full
  50-call BO budget (right), for six cross-task transfer interfaces: a
  single-task baseline with no transfer (STBO), two cross-task interfaces
  that leave initialization uniform and vary only the surrogate/acquisition
  (MTBO, a shared GP; POGPE/SGPE, product-of-experts over frozen per-task
  GPs), and two interfaces that instead learn a task-conditioned
  initialization (BOLT, supervised imitation of high-performing standalone
  candidates; ORPT, trained on downstream rollout outcomes at horizon
  $H=1$). STBO has no training-task dependence and is plotted as a
  reference. The $y$-axis is logarithmic: BOLT, ORPT, MTBO, and STBO
  converge to within 2-3 orders of magnitude of each other while POGPE and
  SGPE remain near their initial regret, a range a linear axis cannot show
  without compressing the former group into visual indistinguishability.}
  \label{fig:motivation-a-scaling}
\end{figure}

\begin{figure}[t]
  \centering
  \includegraphics[width=0.48\textwidth]{figures/motivation_A_trajectory_T50.pdf}
  \caption{Simple regret to the verified global optimum across the full
  50-call BO budget, for all six methods, at the final training checkpoint
  ($T=50$ completed training tasks). Mean over 20 held-out tasks.}
  \label{fig:motivation-a-trajectory}
\end{figure}

% ============================================================
% Experiment B: does standalone quality agree with downstream utility?
% ============================================================

\begin{figure}[t]
  \centering
  \includegraphics[width=0.48\textwidth]{figures/motivation_B_reversal.pdf}
  \caption{Rate at which the sign of the standalone objective-value
  difference between two candidates ($\Delta_0$) disagrees with the sign of
  their matched-intervention downstream-utility difference ($\Delta_h$,
  averaged over $M=3$ shared background pools) at rollout horizon $h$, for
  450 candidate pairs drawn from 10 training-task trajectories (10
  candidates/task). Shown with and without the $z$-score reliability
  filter used by the real preference-construction pipeline. The dotted line
  marks $h=1$, the horizon the main synthetic ORPT experiments train on.
  Near-tied $\Delta_0$ or $\Delta_h$ values are excluded from the rate
  rather than counted toward either outcome. Error bars are 95\% normal-
  approximation confidence intervals.}
  \label{fig:motivation-b-reversal}
\end{figure}

\begin{figure}[t]
  \centering
  \includegraphics[width=0.48\textwidth]{figures/motivation_B_horizon_agreement.pdf}
  \caption{Rate at which the sign of $\Delta_{h}$ at evaluation horizon $h$
  agrees with the sign of $\Delta_{H=1}$ (the training horizon), for the
  same 450 candidate pairs as Figure~\ref{fig:motivation-b-reversal}, with
  and without the reliability filter. Error bars are 95\% normal-
  approximation confidence intervals; the filtered series' sample size
  shrinks at long horizons as more pairs' $\Delta_h$ falls below the
  reliability threshold.}
  \label{fig:motivation-b-agreement}
\end{figure}

% ============================================================
% Experiment C: what should a learned initializer optimize?
% ============================================================

\begin{figure}[t]
  \centering
  \includegraphics[width=0.48\textwidth]{figures/motivation_C_initializer_objective.pdf}
  \caption{Simple regret to the verified global optimum across the 50-call
  BO budget at $T=50$, mean over 20 held-out tasks, for six initializer
  variants that share the same base model, training tasks, held-out tasks,
  proposal budget, downstream BO procedure, and decoding protocol, varying
  only the initializer: task-independent uniform sampling (Random);
  task-independent reuse of the highest-scoring points observed across all
  50 training tasks (Prior-best reuse); a $k$-nearest-neighbor regression
  from scalar task context to each training task's best observed point
  (Context-to-optimum regression); supervised imitation of each task's own
  high-performing candidates (BOLT, i.e.\ Top-K SFT); and two preference-
  tuned variants trained with the same pipeline, differing only in whether
  preferences are labeled by standalone objective value (H=0) or by
  downstream rollout outcome (ORPT, H=1). ORPT (H=3) is shown as a
  secondary variant. Inset: the same curves restricted to oracle calls
  30--40, the region of the main panel where BOLT, the two preference-tuned
  variants, and Prior-best reuse have all converged close enough to zero to
  be indistinguishable at the main panel's scale.}
  \label{fig:motivation-c}
\end{figure}

% Table: see results/figures/motivation_C_table.tex (generated directly by
% figures/figure_c.py::generate_table, not duplicated here to avoid the two
% copies drifting apart).

% Not one of the doc's named figures -- added on request, as a diagnostic
% companion to the trajectory plot above (figures/figure_c.py::
% generate_init_pools). Do not cite as statistical evidence; same status as
% the doc's own optional qualitative figure for Experiment B.
\begin{figure}[t]
  \centering
  \includegraphics[width=\textwidth]{figures/motivation_C_init_pools.pdf}
  \caption{Each initializer's raw $\texttt{init\_size}$ proposed points,
  plotted directly in the 2D Branin input space against that task's own
  contour and numerically verified global optimum, for three held-out tasks
  ($t=0.0, 0.5, 0.95$; first, middle, and last of the 20, fixed and
  deterministic). The contour window for each panel spans the local minimum
  to $+20$ (the well-to-pass height along the valley at $t=0$), so the
  three-well structure and its flattening at high $t$ are both shown at a
  comparable visual scale rather than being compressed by the family's full
  value range. BOLT, both preference-tuned ORPT variants, and Prior-best
  reuse consistently cluster around the same single well across all three
  tasks shown; Context-to-optimum regression's single predicted point is
  visibly outside the valley at $t=0.5$.}
  \label{fig:motivation-c-init-pools}
\end{figure}
